\section{Delivery and control}
\label{sec:delivery}
% Generated by scripts/export_section7.py — do not edit by hand.
\newcommand{\PEPdoaAud}{\$20{,}000\xspace}
\newcommand{\PEPdoaK}{\$20k\xspace}
\newcommand{\PEPdoaHours}{48\xspace}
\newcommand{\PEPcontingency}{\$218{,}700\xspace}
\newcommand{\PEPbase}{\$1{,}500{,}000\xspace}
\newcommand{\PEPbac}{\$1{,}718{,}700\xspace}
\newcommand{\PEPoptionC}{\$21{,}000\xspace}
\newcommand{\PEPpmName}{Abdul Wahab Obeid\xspace}
\newcommand{\PEPmTwoShort}{2~Oct\xspace}
\newcommand{\PEPmFourShort}{15~Oct\xspace}
\newcommand{\PEPmFiveShort}{30~Oct\xspace}
\newcommand{\PEPmSixShort}{10~Nov\xspace}
\newcommand{\PEPmSevenShort}{15~Nov\xspace}
\newcommand{\PEPmFourLong}{15~October\xspace}
\newcommand{\PEPmFiveLong}{30~October\xspace}
\newcommand{\PEPmSevenLong}{15~November\xspace}
\newcommand{\PEPgreenLo}{0.95}
\newcommand{\PEPgreenHi}{1.05}
\newcommand{\PEPamberLo}{0.90}
\newcommand{\PEPredLt}{0.90}
\newcommand{\PEPtZeroSlip}{2\xspace}
\newcommand{\PEPrOiImpact}{\$90{,}000\xspace}
\newcommand{\PEPrOiEmv}{\$27{,}000\xspace}
\newcommand{\PEPrViiImpact}{\$70{,}000\xspace}
\newcommand{\PEPrViiEmv}{\$21{,}000\xspace}
\newcommand{\PEPpvmFour}{\$865{,}578\xspace}
\newcommand{\PEPpvmFournum}{865{,}578}
\newcommand{\PEPbasenum}{1{,}500{,}000}
\newcommand{\PEPbacnum}{1{,}718{,}700}
\newcommand{\PEPonPlanEV}{\$865{,}578\xspace}
\newcommand{\PEPonPlanEVnum}{865{,}578}
\newcommand{\PEPonPlanAC}{\$901{,}000\xspace}
\newcommand{\PEPonPlanACnum}{901{,}000}
\newcommand{\PEPonPlanCPI}{0.96}
\newcommand{\PEPonPlanSPI}{1.00}
\newcommand{\PEPonPlanEac}{\$1{,}561{,}856\xspace}
\newcommand{\PEPonPlanEacnum}{1{,}561{,}856}
\newcommand{\PEPonPlanIeac}{\$1{,}790{,}313\xspace}
\newcommand{\PEPonPlanIeacnum}{1{,}790{,}313}
\newcommand{\PEPlateEV}{\$800{,}178\xspace}
\newcommand{\PEPlateEVnum}{800{,}178}
\newcommand{\PEPlateAC}{\$870{,}000\xspace}
\newcommand{\PEPlateACnum}{870{,}000}
\newcommand{\PEPlateCPI}{0.92}
\newcommand{\PEPlateSPI}{0.92}
\newcommand{\PEPlateEac}{\$1{,}630{,}676\xspace}
\newcommand{\PEPlateEacnum}{1{,}630{,}676}
\newcommand{\PEPmFourGap}{\$0\xspace}
\newcommand{\PEPpmBox}{%
\textbf{Campaign PM}\\[-0.05em]\textbf{\mbox{\PEPpmName}}\\[0.12em]
  \scriptsize DoA: $\le$ \PEPdoaK{} \emph{and} $\le$ \PEPdoaHours{}~h\\[-0.05em]
  \scriptsize Cannot move M-5 (\PEPmFiveShort) or M-7 (\PEPmSevenShort)}


This section turns the baselines in Sections~\ref{sec:schedule}--\ref{sec:risk} into decision rights, procurement, hold points, Bowen HSE, and the earned-value framework read at M-4. The campaign has not been executed.

\subsection{Governance and reporting}
\label{sec:delivery-governance}

Governance is the organisation that will stop work. The Gilmour Executive Board is the sponsor. Abdul Wahab Obeid is Campaign Project Manager. Figure~\ref{fig:governance} and Table~\ref{tab:raci} name the consultancy leads and the field titles.

Delegation is numeric.
\begin{itemize}[leftmargin=1.4em,itemsep=0.25em]
\item Abdul may approve a variance of at most \PEPdoaAud \emph{and} at most \PEPdoaHours~hours, and only if M-5 (\PEPmFiveLong, A-133) and M-7 (\PEPmSevenLong, A-143) do not move.
\item Anything else---a contingency draw, Option~C's \PEPoptionC, a hold-point waiver, a $T$-0 change---goes to the Board within 24~hours.
\end{itemize}
R-16 (RSO versus Launch Director deadlock) escalates 24~hours before A-141's late finish of \PEPmSixShort, not inside the review room. The \PEPdoaAud ceiling is below the dollar \emph{impact} of the hot-fire rows that would recycle the pad (R-01 \PEPrOiImpact; R-07 \PEPrViiImpact). Residual EMVs on those rows are \PEPrOiEmv and \PEPrViiEmv because treatments already sit in the work baseline; a genuine recycle still cannot be buried in PM discretion.

Cadence matches the physical work.
\begin{itemize}[leftmargin=1.4em,itemsep=0.25em]
\item Field leads hold a daily pad standup.
\item After Gate~1 (A-111, \PEPmTwoShort) Ilsa publishes a weekly EVM dashboard from A-112.
\item Abdul takes a fortnightly pack to the sponsor.
\end{itemize}
Daily standup does not rewrite Table~\ref{tab:precedence}. A-112 has 53~days of float so the control system does not compete with stacking for pad occupancy.

\begin{figure}[H]
\centering
\resizebox{\textwidth}{!}{%
\begin{tikzpicture}[
  font=\sffamily,
  >=Stealth,
  board/.style={draw=navy, fill=navy, text=white, rounded corners=2pt,
    inner xsep=8pt, inner ysep=6pt, font=\bfseries\small, align=center, text width=8.4cm, outer sep=1.5pt},
  pm/.style={draw=navy, fill=teal, text=white, rounded corners=2pt,
    inner xsep=6pt, inner ysep=5pt, font=\small, align=center, text width=7.6cm, outer sep=1.5pt},
  lead/.style={draw=navy, fill=white, rounded corners=2pt,
    inner xsep=4pt, inner ysep=4pt, font=\scriptsize, align=center,
    text width=3.85cm, minimum height=1.22cm, outer sep=1.5pt},
  field/.style={draw=navy!70, fill=white, rounded corners=1pt,
    inner xsep=2.5pt, inner ysep=3pt, font=\scriptsize, align=center,
    text width=2.28cm, minimum height=1.08cm, outer sep=1.5pt},
  note/.style={draw=crimson, dashed, thick, rounded corners=2pt, fill=white,
    inner xsep=6pt, inner ysep=4pt, font=\scriptsize, align=center, text width=15.4cm}
]
\node[board] (board) {Gilmour Executive Board \\[0.12em] \normalfont\scriptsize Sponsor --- contingency, M-5 / M-7, hold-point waivers};
\node[pm, below=7mm of board] (pm) {\PEPpmBox};
\draw[navy, thick, ->, shorten >=2pt] (board.south) -- (pm.north);

\node[lead, anchor=north] (hadi)  at (-4.55,-4.55) {\textbf{\mbox{Hadi Alkhub}}\\Lead Planner\\schedule, network, compression};
\node[lead, anchor=north] (ziyad) at ( 0.00,-4.55) {\textbf{\mbox{Ziyad Alahmari}}\\Risk and Quality Lead\\register, hold points, HSE};
\node[lead, anchor=north] (ilsa)  at ( 4.55,-4.55) {\textbf{\mbox{Ilsa Hashmi}}\\Cost and Resource Lead\\PMB, EVM, Juru commercial};
\coordinate (pmfork) at ($(pm.south)!0.55!(ziyad.north)$);
\draw[navy, thick] (pm.south) -- (pmfork);
\foreach \n in {hadi,ziyad,ilsa}
  \draw[navy, thick, ->, shorten >=2.5pt] (pmfork) -| (\n.north);

\node[field, anchor=north] (ld)   at (-6.35,-7.15) {\textbf{Launch}\\Director};
\node[field, anchor=north] (rso)  at (-3.81,-7.15) {\textbf{Range Safety}\\Officer};
\node[field, anchor=north] (gse)  at (-1.27,-7.15) {\textbf{GSE}\\Lead};
\node[field, anchor=north] (avio) at ( 1.27,-7.15) {\textbf{Avionics}\\Specialist};
\node[field, anchor=north] (juru) at ( 3.81,-7.15) {\textbf{Juru Cultural}\\Officer};
\node[field, anchor=north] (dc)   at ( 6.35,-7.15) {\textbf{D\&C package}\\manager};
\coordinate (fieldbar) at ([yshift=-5.5mm]ziyad.south);
\draw[navy, thick] (hadi.south) -- (hadi.south |- fieldbar);
\draw[navy, thick] (ziyad.south) -- (fieldbar);
\draw[navy, thick] (ilsa.south) -- (ilsa.south |- fieldbar);
\draw[navy, thick] (ld.north |- fieldbar) -- (dc.north |- fieldbar);
\foreach \n in {ld,rso,gse,avio,juru,dc}
  \draw[navy!80, ->, shorten >=2.5pt] (\n.north |- fieldbar) -- (\n.north);

\node[note, anchor=north] (cadence) at (0,-8.85) {\textbf{Reporting cadence.} Daily pad standup \quad$\cdot$\quad weekly EVM dashboard (Ilsa) \quad$\cdot$\quad fortnightly sponsor review \quad$\cdot$\quad 24~h to Board if DoA exceeded or R-16 fires.};
\end{tikzpicture}%
}
\caption[Campaign governance structure]{Campaign governance for Bowen pad operations. Named roles, a numeric DoA on the PM, and the reporting cadence. Field titles match Table~\ref{tab:raci}. The Board, not the PM, moves M-5 or M-7 or draws the \PEPcontingency reserve.}
\label{fig:governance}
\end{figure}

\subsection{Delivery model and procurement}
\label{sec:delivery-model}

The chosen delivery model is a managing-contractor / design-and-construct hybrid. Gilmour retains propulsion intellectual property, vehicle integration and the Australian launch licence. Civil works and cryogenic GSE packages are let as design-and-construct, with an interface-control freeze and liquidated damages. That split is the contractual answer to R-07 and R-08: an ICD lag is an LD event on the contractor, and a latent pad-civil defect stays with the D\&C package after HP-1 NDT (R-08 EMV \$0 to Gilmour). Abdul stays Accountable on that row; the package manager stays Responsible.

Turnkey EPC was rejected. After TestFlight~1, giving the propulsion stack and the launch authorisation to one EPC would put flight-critical intellectual property and statutory accountability in a contractor who does not hold the licence \citep{gilmour2025,asa2025}. Two tests decide the model. Scope certainty is high on site operations (seventeen Level-3 packages, locked dates, a \PEPbase\ work estimate) and low on the vehicle. Client capability is the other test: Gilmour can integrate a launch vehicle; pouring pad grout or running a 100~t crawler is not its core work. The charter set the weights at 40/30/30, and A1 did not score the options. This PEP scores the site-ops choice only:
\begin{itemize}[leftmargin=1.4em,itemsep=0.25em]
\item Strategic Fit 40\% (EPC 2/5, D\&C 5/5).
\item Technical Feasibility 30\% (3/5 vs 4/5).
\item Schedule 30\% (3/5 vs 4/5).
\end{itemize}
Weighted 2.6 versus 4.4, so D\&C wins. Alliance contracting was the second rejected alternative: a four-month campaign cannot carry pain-share overhead, and GBRMPA and CASA duties are not commercially shareable \citep{gbrmpa2019,casa2024}.

\subsection{Quality}
\label{sec:delivery-quality}

Quality uses two standards together \citep{sae2016,iso9001}. ISO~9001:2015 is the general rule: plan the work, name who checks it, and do not pass a job that has not been verified. AS9100D adds hold points, extra checks on welds, leaks and ignition, and a formal release before flight. This campaign uses those rules as pad stops.

ISO~9001 maps onto named owners. The PEP is the quality plan and Ziyad is the release authority. Table~\ref{tab:raci} names who is Accountable on each package. Design defects found on the pad go back to R\&D. If HP-1 is unsigned, stacking earns \$0.

AS9100D maps onto the four hold points in Table~\ref{tab:holdpoints}. Work does not pass a hold without Ziyad's release or a Board waiver, which sits outside Abdul's authority. Option~A was rejected because signing HP-4 on incomplete post-fire data would be a hold-point failure.

\begin{table}[H]
\centering
\small
\caption[Quality hold points]{Quality hold points for the Bowen campaign. Each hold is a named artefact, not a meeting. Measurement rules in Section~\ref{sec:delivery-evm} follow the same IDs.}
\label{tab:holdpoints}
\begin{tabularx}{\textwidth}{@{}l l l l X@{}}
\toprule
Hold & Activity & Risk & Artefact & Release \\
\midrule
HP-1 & A-123 (15~Oct) & R-08 & Pad tie-down weld NDT & D\&C mgr / Ziyad \\
HP-2 & A-124 (12~Oct) & R-05 & GSE helium leak test & GSE Lead / Ziyad \\
HP-3 & A-133 (30~Oct) & R-01 & Pre-static-fire ignition interlock & GSE Lead / Launch Director / Ziyad \\
HP-4 & A-141 (10~Nov) & R-16 & LRR sign-off (payload sponsor witnesses) & Launch Director / Abdul / Ziyad / payload sponsor \\
\bottomrule
\end{tabularx}
\end{table}


\begin{itemize}[leftmargin=1.4em,itemsep=0.25em]
\item Discrete tests and flight artefacts---A-132, A-133, A-141, A-143---are 0/100.
\item A-113 is elapsed-day percent-complete because the work is a 76-day wait, not a binary test.
\item Installations---A-123 and A-124---are milestone percent-complete tied to HP-1 and HP-2.
\end{itemize}
LRR waits on A-134 rather than on A-133: HP-4 cannot be signed on incomplete post-fire data, which is why Option~A was rejected \citep{sae2016}.

\subsection{Health, safety and environment}
\label{sec:delivery-hse}

HSE is this pad under the \textit{Work Health and Safety Act 2011} (Qld) \citep{whs2011}.
\begin{itemize}[leftmargin=1.4em,itemsep=0.25em]
\item Liquid oxygen is the process hazard: LOX PPE, oxygen-clean tools, and exclusion during load and detank on A-132 and A-133.
\item This PEP's pad controls---not A1 charter numbers---are a 1.5~km blast zone and twelve simultaneous heads, the same cap on Figure~\ref{fig:resources}.
\item Acoustic treatment is a water deluge so static-fire level at Bowen township stays at or below 135~dB (R-17), with community notice fourteen days prior.
\item Juru artefact buffers and stop-work authority (R-04) are HSE controls: A-114 feeds A-132 as a cultural hold.
\item Finish-line HSE remains as chartered: zero lost-time injuries, zero GBRMPA breaches, and zero Juru heritage breaches.
\item Marine-mammal protocol on A-142 (R-15) is owned by the RSO \citep{gbrmpa2019}.
\end{itemize}

\subsection{Earned value management}
\label{sec:delivery-evm}

Earned value tells the Board whether the \PEPbac\ authorisation is still intact \citep{pmi2019evm}. It is a framework. The campaign has not started, so there is no EV/AC history to plot.

Distributed planned value is \PEPbase, the same total as Table~\ref{tab:wbs-cost} and Figure~\ref{fig:scurve}. Contingency of \PEPcontingency\ is undistributed: not December cash and not monthly PV. Authorised $BAC$ is \PEPbac. Option~C's \PEPoptionC\ sits outside the work baseline; if the Board invokes it, the draw is actual cost charged to contingency.

$SV$ and $CV$ are currency, never days:
\begin{equation}
SV = EV - PV, \qquad CV = EV - AC, \qquad SPI = EV/PV, \qquad CPI = EV/AC.
\label{eq:evm}
\end{equation}
In-flight estimate at completion on the work baseline is
\begin{equation}
EAC = AC + (BAC_{\text{work}} - EV)/CPI,
\label{eq:eac}
\end{equation}
with $BAC_{\text{work}} = \PEPbase$. The same CPI is applied to the authorisation ceiling as $BAC \times AC / EV$, reported against \PEPbac \citep{pmi2019evm}. Finish-line rule: final $AC \le \PEPbac$. A Green CPI of \PEPonPlanCPI\ in October is allowed; finishing above $BAC$ is not.

Thresholds that trigger action, used only in this section, are:

\begin{itemize}[leftmargin=1.4em,itemsep=0.25em]
\item \textbf{Green} --- $\PEPgreenLo \le CPI,SPI \le \PEPgreenHi$: dashboard only.
\item \textbf{Amber} --- $\PEPamberLo \le$ index $< \PEPgreenLo$: Exception Report and a seven-day recovery plan within 48~hours.
\item \textbf{Red} --- index $< \PEPredLt$, \emph{or} $T$-0 slip $> \PEPtZeroSlip$~days: contingency freeze and Board intervention. $SV$ is not written in days.
\end{itemize}

The only dashboard numbers in this PEP are the two illustrative readings at M-4 (\PEPmFourLong~2026) in Table~\ref{tab:evm-m4} and Figure~\ref{fig:evm}. On-plan is Green ($CPI = \PEPonPlanCPI$, $SPI = \PEPonPlanSPI$), but the estimate at completion still sits above \PEPbac, so the finish-line rule is not yet met. The late-stack case is \PEPlateBand\ ($CPI = \PEPlateCPI$, $SPI = \PEPlateSPI$) because HP-1 is unsigned and the \$218{,}000 stacking account earns \$0. Ilsa issues the Exception Report within 48~hours and freezes contingency. If R-14 has fired, the recovery invokes Option~C and crashes A-134 only, not A-141 or A-142.

\begin{table}[H]
\centering
\footnotesize
\setlength{\tabcolsep}{4.0pt}
\caption[Illustrative earned-value status at M-4]{Illustrative earned-value status at M-4 (15~October~2026). The campaign has not been executed. PV is the locked work-PMB value at stacking complete. Indices use the reported two-decimal CPI/SPI from $EV/AC$ and $EV/PV$.}
\label{tab:evm-m4}
\begin{tabularx}{\textwidth}{@{}>{\raggedright\arraybackslash}p{4.55cm} r r r c c X@{}}
\toprule
Status at M-4 & PV & EV & AC & CPI & SPI & Action \\
\midrule
HP-1 signed (A-123 100\%), HP-2 signed, A-113 elapsed-day percent-complete & 865{,}578 & --- & --- & --- & --- & --- \\
Illustrative on-plan & 865{,}578 & 865{,}578 & 901{,}000 & 0.96 & 1.00 & Green --- weekly dashboard \\
Illustrative late stack & 865{,}578 & 647{,}578 & 870{,}000 & 0.74 & 0.75 & Red --- 48~h exception report; contingency freeze \\
\bottomrule
\end{tabularx}
\end{table}


\begin{figure}[!htbp]
\centering
\includegraphics[width=0.98\linewidth]{evm_example.pdf}
\caption[Illustrative earned-value chart at M-4]{Illustrative status at M-4; campaign not yet executed. Upper panel: PV, EV and AC at the single status date \PEPmFourLong~2026 for the on-plan and late-stack cases. Lower panel: Green / Amber / Red CPI--SPI bands, with the on-plan pair (\PEPonPlanCPI, \PEPonPlanSPI) in Green and the late-stack pair (\PEPlateCPI, \PEPlateSPI) in \PEPlateBand. Red is also triggered by $T$-0 slip greater than \PEPtZeroSlip days. No weekly EV/AC history is implied.}
\label{fig:evm}
\end{figure}
